% STATUS: written, 20 October 2026.
%
% Every citation below points at a source I have read. Before submission,
% fill the entries marked VERIFY in references.bib with page numbers and
% DOIs from the publisher, and add any further related work in the three
% groups the section already sets up. Do not add a citation you have not
% opened.

\section{Related Work}

Three bodies of work meet in this problem: the classical treatment of
fragmentation and allocation as an optimization to be solved, recent systems
that adapt placement to a moving workload, and the use of learning to replace
hand-written database administration.

\subsection{Fragmentation and allocation as an optimization problem}

The question of how to divide a database and where to put the pieces is as old
as distributed databases themselves, and the classical treatment is to state it
as an optimization problem and solve it once. A cost model scores a candidate
arrangement against an expected workload, and a solver searches for the
arrangement that minimizes it. The result is then installed and the database
runs with it.

The weakness is in the word once. A solved-once answer is correct for the
workload it was solved against and degrades as the real workload departs from
it. Where the departure is slow this is tolerable, and where the workload has a
geographic centre that moves on a daily cycle it is not. That observation
motivates the systems in the next group, which re-solve the problem as the
workload changes.

% TODO: cite two or three classical treatments here. Candidates to read:
% vertical and horizontal partitioning in distributed database design, and
% the allocation problem as stated in the standard textbook treatment.

\subsection{Adapting placement to a moving workload}

Lion~\cite{lion2024} is the closest system to this work in intent. It observes
that a transaction is slow mainly because it spans machines, and it therefore
arranges replicas so that some single node holds a copy of every partition a
given transaction touches, allowing that transaction to commit locally rather
than through the multi-round agreement protocol that a distributed commit
requires. A graph-based analyzer groups partitions that are accessed together,
a cost model assigns those groups to nodes subject to a load-balance
constraint, and an LSTM predictor forecasts the workload so that replicas are
provisioned before the shift rather than after it. Lion reports up to 2.7 times
higher throughput and 76.4\% better scalability than the approaches it is
compared against.

Lion establishes three things this work depends on. Adapting placement to a
moving workload is worth a large amount of real throughput. Migration cost and
load balance both belong in the objective rather than latency alone.
Anticipating a shift beats reacting to one, by roughly a factor of four in
adaptation speed.

It also leaves three openings. The adaptation is a heuristic pipeline rather
than a learned policy: the LSTM forecasts, but a hand-written cost model
decides, so the system never improves from the outcomes of its own decisions.
Geographic distance does not appear in its cost function, which is expressed in
bytes moved, so two placements that differ by an ocean score identically.
Fragment boundaries are fixed, so Lion moves replicas of a partition but never
splits one.

GeoTP~\cite{geotp2025} supplies the geographic dimension that Lion does not
have. It targets database middleware connecting heterogeneous data sources
across regions, where the dominant costs are the network latency between the
middleware and its data sources and the length of time a transaction holds
locks whilst waiting on remote work. GeoTP reduces both: a decentralized
prepare phase cuts the number of network round trips a distributed transaction
needs, and a latency-aware scheduler delays lock acquisition so that the
contention window is shorter. It reports up to 17.7 times improvement on YCSB
and TPC-C.

GeoTP shows that latency awareness is worth building a scheduler around, and
its cost accounting is the model for the one used here. It optimizes the
execution of transactions against a given placement, however, and does not
change where the data sits. The two papers are therefore complementary, and
the gap between them is where this work sits: Lion adapts placement without
seeing geography, GeoTP sees geography without adapting placement.

\subsection{Learning to administer a database}

A separate line of work replaces hand-written database administration with
learned policies, in index selection, in configuration tuning, and in query
optimization. The common argument is that these decisions are made repeatedly
against a changing workload, that the quality of a decision is measurable after
the fact, and that a system able to learn from its own outcomes should
eventually beat a fixed rule written by a person who could not anticipate every
workload.

Placement is a natural member of that family and has received less attention
than the others, which is the opportunity this work takes up.

% TODO: cite representative learned index selection, knob tuning and query
% optimization work. Read at least one from each before citing it.

\subsection{What is new here}

This work learns fragmentation and geographic placement \emph{jointly}, as one
policy, for transactional workloads. Lion learns neither and adapts only
replica placement. GeoTP optimizes execution rather than placement. The learned
administration literature addresses other decisions. Stating the contribution
precisely: the agent chooses, in one action space, where a shard's primary copy
sits, where its replicas sit, and where the boundary between two shards falls,
and it is rewarded on a cost that combines transaction latency, the volume of
data moved, and load imbalance.

I state one limitation of the contribution up front, because the evaluation
reports it. The agent described here is wired, trained and measured, and it
does not yet outperform the hand-written heuristic it is meant to replace.
What this paper establishes is the problem, the measurement apparatus, the
baselines, and the conditions under which a learned policy has room to win.
